\renewcommand{\Anno}{2005}

\renewcommand{\NumeroProblema}{4}

\renewcommand{\Testo}{Una commissione deve assegnare una borsa di
studio scegliendo tra due candidati: Paolo e Francesca.  Il concorso
consiste in una prova scritta che comprende $n$ problemi.  Ogni
problema pu\`{o} essere valutato ``risolto'', oppure ``non risolto'',
senza gradazioni intermedie.  Il candidato che risolve il maggior
numero di problemi ottiene la borsa di studio; in caso di parit\`{a}
la commissione assegna la borsa ``tirando la monetina'', cio\`{e}
scegliendo con uguale probabilit\`{a} tra i due candidati.  Il giorno
del concorso la preparazione di Paolo gli permette di risolvere ogni
problema con probabilit\`{a} 1/2, mentre quella di Francesca gli
consente di risolvere ogni problema con probabilit\`{a} 2/3.
	\begin{enumerate}
		\item  Nel caso $n=1$, determinare la probabilit\`{a} che Paolo 
		vinca la borsa di studio.
	
		\item Nel caso $n=2$, determinare la probabilit\`{a} che
		Francesca vinca la borsa di studio.
	
		\item  Dimostrare che, nel caso in cui vengano assegnati $n\geq 2$ 
		problemi, la commissione dovr\`{a} ricorrere alla monetina con 
		probabilit\`{a} strettamente maggiore di $(1/2)^{n}$.
\end{enumerate} }

\renewcommand{\Risposta}{\begin{enumerate}
		\item  5/12
	
		\item  5/8
	
		\item  \ldots
	\end{enumerate}}

\renewcommand{\Aiutino}{\begin{enumerate}
	\renewcommand{\labelenumi}{(\alph{enumi})} 
	\item  Calcolare la probabilit\`{a} dei 4 possibili punteggi finali. 
	Sommare le probabilit\`{a} di quelli favorevoli a Paolo pi\`{u} 
	met\`{a} dei pareggi.

	\item Calcolare, per ciascuno dei 2 concorrenti, la
	probabilit\`{a} che risolva 0, 1 o 2 esercizi.  Calcolare quindi
	la probabilit\`{a} dei 9 possibili punteggi finali.  Sommare le
	probabilit\`{a} di quelli favorevoli a Francesca pi\`{u} met\`{a} dei
	pareggi.

	\item  Vi sono almeno due approcci.
	\begin{itemize}
		\item  Calcolare la probabilit\`{a} che Paolo e Francesca pareggino 
		ottenendo lo stesso punteggio in \emph{ogni} esercizio.
	
		\item  Esprimere la probabilit\`{a} di pareggio finale come una 
		sommatoria che coinvolge coefficienti binomiali al quadrato. 
		Minorare tale sommatoria togliendo i quadrati ai binomiali e 
		concludere con il binomio di Newton.
	\end{itemize}
\end{enumerate}}

\renewcommand{\SoluzioneUfficiale}{\TitoloSoluzione\DomandaA I 
possibili punteggi finali sono 4, a seconda che Paolo e Francesca 
risolvano 0 o 1 esercizio. Le probabilit\`{a} di questi 4 eventi sono 
ricavate nella seguente tabella.
\begin{center}
	\renewcommand{\arraystretch}{1.3}
	\begin{tabular}{|c|c||c|c|}
		\cline{3-4}
		\multicolumn{2}{c|}{} & Francesca 0 & Francesca 1 \\
		\cline{3-4}
		\multicolumn{2}{c|}{} & 1/3 & 2/3 \\
		\cline{3-4}\noalign{\vspace{0.4ex}}\hline
		Paolo 0 & 1/2 & 1/6 & 1/3  \\
		\hline
		Paolo 1 & 1/2 & 1/6 & 1/3  \\
		\hline
	\end{tabular}
\end{center}

La probabilit\`{a} di vittoria di Paolo \`{e} dunque data dalla somma 
della probabilit\`{a} che Paolo vinca 1--0, pi\`{u} met\`{a} della 
probabilit\`{a} dei pareggi (0--0 e 1--1) ed \`{e} pertanto
$$\frac{1}{6}+\frac{1}{2}\left(\frac{1}{6}+\frac{1}{3}\right)=
\frac{1}{6}+\frac{1}{4}=\frac{5}{12}.$$

\DomandaB Ora sia Paolo sia Francesca possono risolvere 0, 1 o 2 
esercizi, e pertanto gli eventi possibili sono 9. Paolo risolver\`{a} 
0 problemi con probabilit\`{a} $1/2\cdot 1/2=1/4$, 2 problemi con 
probabilit\`{a} $1/2\cdot 1/2=1/4$ e, per differenza, 1 problema con 
probabilit\`{a} 1/2. Francesca risolver\`{a} 
0 problemi con probabilit\`{a} $1/3\cdot 1/3=1/9$, 2 problemi con 
probabilit\`{a} $2/3\cdot 2/3=4/9$ e, per differenza, 1 problema con 
probabilit\`{a} 4/9.

Ne segue che le probabilit\`{a} dei 9 possibili punteggi finali sono 
quelle ricavate nella seguente tabella
\begin{center}
	\renewcommand{\arraystretch}{1.3}
	\begin{tabular}{|c|c||c|c|c|}
		\cline{3-5}
		\multicolumn{2}{c|}{} & Francesca 0 & Francesca 1 & Francesca 2\\
		\cline{3-5}
		\multicolumn{2}{c|}{} & 1/9 & 4/9 & 4/9 \\
		\cline{3-5}\noalign{\vspace{0.4ex}}\hline
		Paolo 0 & 1/4 & 1/36 & 1/9 & 1/9 \\
		\hline
		Paolo 1 & 1/2 & 1/18 & 2/9 & 2/9\\
		\hline
		Paolo 2 & 1/4 & 1/36 & 1/9 & 1/9\\
		\hline
	\end{tabular}
\end{center}

La probabilit\`{a} di vittoria di Francesca \`{e} dunque data dalla
somma delle probabilit\`{a} dei 3 eventi in cui Francesca totalizza
pi\`{u} punti di Paolo, pi\`{u} met\`{a} delle probabilit\`{a} dei 3
casi di pareggio, ed \`{e} pertanto
$$\frac{1}{9}+\frac{1}{9}+\frac{2}{9}+
\frac{1}{2}\left(\frac{1}{36}+\frac{2}{9}+\frac{1}{9}\right)=
\frac{4}{9}+\frac{13}{72}=\frac{45}{72}=\frac{5}{8}.$$

\DomandaC Determiniamo la probabilit\`{a} che Paolo e Francesca
pareggino ottenendo lo stesso punteggio in tutti gli esercizi.
Poich\'{e} Paolo ha probabilit\`{a} 1/2 sia di risolvere, sia di non
risolvere un dato esercizio, egli ha in particolare probabilit\`{a}
1/2 di ottenere, in quell'esercizio, lo stesso risultato di Francesca.
Ne segue che la probabilit\`{a} che Paolo ottenga lo stesso risultato
di Francesca in tutti gli $n$ problemi \`{e} $1/2^{n}$, e questo
dimostra che per ogni intero positivo $n$ la probabilit\`{a} di
pareggio \`{e} maggiore od uguale di $1/2^{n}$.  Se $n\geq 2$,
potendosi avere il pareggio anche con punteggi parziali diversi, la
disuguaglianza diventa stretta.}

\renewcommand{\SoluzioniAlternative}{\TitoloSoluzioneAlternativaA La 
probabilit\`{a} che Paolo risolva $k$ esercizi su $n$ \`{e}
$$\binom{n}{k}\frac{1}{2^{n}},$$
la probabilit\`{a} che Francesca risolva $k$ esercizi su $n$ \`{e}
$$\binom{n}{k}\left(\frac{2}{3}\right)^{k}\frac{1}{3^{n-k}}=
\binom{n}{k}\frac{2^{k}}{3^{n}}.$$

La giustificazione di tali formule \`{e} piuttosto classica: per 
Francesca, per esempio, il binomiale rappresenta i modi di scegliere 
$k$ esercizi da risolvere su $n$, il fattore $(2/3)^{k}$ rappresenta 
la sua probabilit\`{a} di risolvere $k$ esercizi fissati e 
$(1/3)^{n-k}$ la sua probabilit\`{a} di non risolvere i rimanenti 
$n-k$.

Ne segue che la probabilit\`{a} di pareggio \`{e} data da
$$\sum_{k=0}^{n}\binom{n}{k}^{2}\frac{1}{2^{n}}\cdot\frac{2^{k}}{3^{n}}=
\frac{1}{6^{n}}\sum_{k=0}^{n}\binom{n}{k}^{2}2^{k}.$$

Possiamo ora maggiorare la sommatoria togliendo il quadrato ai 
coefficienti binomiali e notando che per $n\geq 2$ la maggiorazione 
\`{e} stretta in quanto almeno uno dei coefficienti binomiali \`{e} 
strettamente maggiore di 1. Dalla formula del binomio di Newton 
otteniamo quindi che
$$\frac{1}{6^{n}}\sum_{k=0}^{n}\binom{n}{k}^{2}2^{k}>
\frac{1}{6^{n}}\sum_{k=0}^{n}\binom{n}{k}2^{k}\cdot 1^{n-k}=
\frac{1}{6^{n}}(2+1)^{n}=\frac{1}{2^{n}},$$
che \`{e} il risultato cercato.}

\renewcommand{\Osservazioni}{\TitoloOsservazione Nel linguaggio della 
teoria della probabilit\`{a}, le prime due formule della soluzione 
alternativa dicono semplicemente che la somma di $n$ variabili 
aleatorie di Bernoulli indipendenti \`{e} una variabile aleatoria 
binomiale.
}

\renewcommand{\Commenti}{
}